\section{Post-Recall Evidence Evaluation}
\label{sec:framework}

We study memory-based question answering after the relevant evidence has been recalled.
The point is not that retrieval is unimportant: retrieval determines whether the needed evidence is available.
Our focus is the next stage.
Once evidence is available, the system must expose a decision-sufficient view, bind the relevant variables, execute the required operations, and produce an answer whose denotation is valid under the evidence.
We refer to retrieval-side mechanisms as \emph{access geometry} and post-recall reasoning mechanisms as \emph{evaluation algebra}.

\subsection{Evidence-Defined Validity}
\label{subsec:evidence-validity}

Let memory be a finite typed evidence structure $\mathcal A=(E,R_1,\ldots,R_m)$, where $E$ is a set of evidence items and the relations $R_j$ range over items, entities, timestamps, sources, amounts, statuses, locations, events, and states.
Relations may include attributes such as $\mathsf{amount}(e,a)$, $\mathsf{status}(e,s)$, and $\mathsf{time}(e,t)$, as well as structural edges such as $\mathsf{newer\text{-}than}(e_i,e_j)$, $\mathsf{overrides}(e_i,e_j)$, $\mathsf{same\text{-}event}(e_i,e_j)$, and $\mathsf{contradicts}(e_i,e_j)$.

A query $q$ induces a semantic answer relation $\varphi_q$.
The evidence-defined answer set is
\begin{equation}
\mathcal Z_q(\mathcal A)=\{z\in\mathcal Z:\mathcal A\models\varphi_q(z)\},
\label{eq:semantic-answer-set}
\end{equation}
where $\mathcal Z$ may contain entities, numbers, records, lists, state decisions, claim sets, or structured tuples.
A natural-language response $y\in\Sigma^\ast$ is evaluated through a partial semantic projection $\eta:\Sigma^\ast\rightharpoonup\mathcal Z$.
The answer is evidence-valid iff
\begin{equation}
\eta(y)\in\mathcal Z_q(\mathcal A).
\label{eq:evidence-valid}
\end{equation}
This target is independent of surface fluency: the model is correct only when its output denotes an element of the evidence-defined answer set.

Let an access module return a query-conditioned substructure $\mathcal A_q=R_{\mathsf{access}}(q,\mathcal A)\preceq\mathcal A$.
Let $\mathcal A_q^\star$ denote a sufficient evidence substructure for answering $q$.
We call access oracle when $\mathcal A_q^\star\preceq \mathcal A_q$.
Post-recall evaluation studies the conditional regime in which oracle evidence is present.
Successful answering then requires
\begin{equation}
\mathsf{Success}
=
\mathsf{AccessOK}
\land
\mathsf{ViewOK}
\land
\mathsf{EvalOK}.
\label{eq:success-factor}
\end{equation}
Our experiments condition on $\mathsf{AccessOK}$ and measure the remaining failures due to view insufficiency and evaluation error.

\subsection{Decision-Sufficient Views}
\label{subsec:decision-sufficient-views}

The recalled evidence is not consumed by the model as $\mathcal A_q$.
It is exposed through a view and a serialization,
\begin{equation}
\mathcal A_q
\xrightarrow{\mathsf{View}_q}
S_q
\xrightarrow{\mathsf{Ser}_q}
V_q .
\label{eq:view-pipeline}
\end{equation}
The view decides which fields, timestamps, statuses, source identifiers, modality markers, and relations are exposed.
The serialization decides whether these elements appear as flat text, source-indexed records, a table, explicit relation triples, state slots, or a tool trace.

A view can include all retrieved evidence and still be insufficient for the decision.

\paragraph{Definition 1 (Decision-sufficient view).}
A view $V_q$ is decision-sufficient for query $q$ if, for any two evidence structures $\mathcal A_1,\mathcal A_2$,
\begin{equation}
V_q(\mathcal A_1)=V_q(\mathcal A_2)
\Longrightarrow
\mathcal Z_q(\mathcal A_1)=\mathcal Z_q(\mathcal A_2).
\label{eq:view-sufficient}
\end{equation}

The contrapositive is the useful part.
If
\begin{equation}
V_q(\mathcal A_1)=V_q(\mathcal A_2),
\qquad
\mathcal Z_q(\mathcal A_1)\cap\mathcal Z_q(\mathcal A_2)=\emptyset,
\label{eq:view-collapse}
\end{equation}
then no decoder observing only $V_q$ can be correct on both worlds.
The decoder receives the same input in two worlds that require disjoint semantic answers.
Thus, answer-relevant distinctions lost by the view are unrecoverable from the view alone.

\paragraph{Proposition 1 (View collapse impossibility).}
If Equation~\ref{eq:view-collapse} holds, no deterministic or randomized decoder that observes only $V_q$ can be evidence-valid on both $\mathcal A_1$ and $\mathcal A_2$ with probability one.
The proof is in Appendix~\ref{app:proof-view-collapse}.

This is the first nontrivial bottleneck after recall.
The problem is not whether the evidence exists somewhere in memory, but whether the exposed view preserves the variables and relations needed by the query.
A flat text view may be decision-sufficient for a single-field lookup, but insufficient for temporal override, source-priority resolution, contradiction handling, or abstention under missingness.
We therefore treat view sufficiency as a query-view property, not as a global property of a format.

Operationally, our benchmark pairs each query with a gold evaluation program.
The program specifies the fields, relations, and state variables required for the answer.
This lets us estimate whether a view is decision-sufficient by checking whether it exposes the required variables in a usable form.

\subsection{Dynamic Propagation of Binding Error}
\label{subsec:dynamic-binding}

Post-recall evaluation often requires binding multiple variables before applying an operator: a receipt to an amount, a booking to a date range, an old record to an update, a claim to its source, or a conflicting item to a priority rule.
Attention provides soft alignment, not discrete binding.
At evaluation step $t$, let the model read a soft-bound value with attention weights $\alpha_{ti}\ge0$ and $\sum_i\alpha_{ti}=1$:
\begin{equation}
\tilde v_t=\sum_i \alpha_{ti}v_i .
\label{eq:soft-bound-value}
\end{equation}
Let $S_t$ be the correct support set for the variable needed at step $t$, suppose all values in $S_t$ share the target value $v_t^\star$, and define the binding mass
\begin{equation}
B_t=\sum_{i\in S_t}\alpha_{ti},
\qquad
D_t=\max_i\lVert v_i-v_t^\star\rVert_2 .
\label{eq:binding-mass}
\end{equation}

The important question is not the static perturbation alone, but how it propagates through a multi-step evaluation trajectory.
Let
\begin{equation}
h_{t+1}=T_t(h_t,\tilde v_t),
\qquad
h_{t+1}^{\star}=T_t(h_t^\star,v_t^\star),
\label{eq:evaluation-dynamics}
\end{equation}
where $h_t$ is the actual decision state and $h_t^\star$ is the ideal state under correct binding.
Define $e_t=\lVert h_t-h_t^\star\rVert_2$.
Assume the local update is Lipschitz in the current state and the bound variable:
\begin{equation}
\lVert T_t(h,v)-T_t(h',v')\rVert_2
\le
\rho_t\lVert h-h'\rVert_2+\beta_t\lVert v-v'\rVert_2 .
\label{eq:transition-lipschitz}
\end{equation}

\paragraph{Proposition 2 (Dynamic propagation of binding error).}
Under the above conditions,
\begin{equation}
e_{t+1}
\le
\rho_t e_t+\beta_t(1-B_t)D_t .
\label{eq:dynamic-one-step}
\end{equation}
Consequently,
\begin{equation}
\begin{aligned}
e_T
\le&
\left(\prod_{t=0}^{T-1}\rho_t\right)e_0\\
&+
\sum_{s=0}^{T-1}
\left(\prod_{u=s+1}^{T-1}\rho_u\right)
\beta_s(1-B_s)D_s .
\end{aligned}
\label{eq:dynamic-unrolled}
\end{equation}
The proof is in Appendix~\ref{app:proof-dynamic-binding}.

This bound is the main theoretical prediction.
Post-recall errors should not merely depend on whether evidence was recalled.
They should depend on the length and stability of the evaluation trajectory.
When $\rho_t<1$, binding errors are damped; when $\rho_t\approx1$, they accumulate; when $\rho_t>1$, early binding errors may be amplified.
Therefore, under oracle evidence, error should increase with binding count, operator count, override depth, conflict degree, and missingness decisions.
This motivates the complexity sweeps in Section~\ref{sec:atm}.

The same bound also explains why structured views and harnesses can help.
A structured view can increase $B_t$ by making the correct support easier to identify.
A decision-state view can reduce $T$ by exposing the operands needed for evaluation.
An operator harness can remove unstable computation from the model trajectory and execute it externally.

\subsection{Decision State and Harnessed Evaluation}
\label{subsec:decision-state-harness}

A decision state is the set of typed variables needed by the gold evaluation program.
For lookup queries, this may be one field.
For temporal, conflict, or aggregation queries, it may include candidate records, stale records, timestamps, source priorities, override edges, conflict flags, missingness indicators, and numeric operands.
The purpose of a decision state is not to preserve all recalled text, but to preserve the variables needed for the answer.

Language models produce proposal distributions over strings, $p_\theta(y\mid q,V_q)$.
Evidence semantics defines which denotations are valid, $\eta(y)\in\mathcal Z_q(\mathcal A)$.
Prompting, chain-of-thought, and scratchpads reshape the proposal distribution.
They may improve the internal trajectory, but they do not enforce evidence validity.

Harnessed evaluation changes the interface.
Operator harnesses execute discrete operations such as $\mathsf{sum}$, $\mathsf{count}$, $\mathsf{latest}$, $\mathsf{override}$, $\mathsf{join}$, and status comparison.
Verifier harnesses reject candidates that violate source, type, timestamp, relation, support, or missingness constraints.
The verifier does not merely prefer correct answers; it deletes invalid ones.
When no proposal survives, the system should abstain, not from uncertainty, but from evidence-level non-authorization.

This distinction guides the empirical design.
Prompting-only methods are tested as proposal-shaping baselines.
Operator and verifier harnesses are tested as evaluation-centric interventions that either externalize unstable operations or restrict admissible outputs.
